\label{sec:results}

\subsection{Open-Seat Counts}

Table~\ref{tab:open-results} presents the strict and broad open-seat counts for
\textsc{bisect} algorithmic maps and enacted 2022 congressional maps in all three states.

\begin{table}[ht]
\centering
\begin{threeparttable}
\caption{Open-seat counts: \textsc{bisect} vs.\ enacted 2022 maps vs.\ baseline.}
\label{tab:open-results}
\begin{tabular}{l r r r r r}
\toprule
 & & \multicolumn{2}{c}{Open seats (strict)} & \multicolumn{2}{c}{Open seats (broad)} \\
\cmidrule(lr){3-4} \cmidrule(lr){5-6}
State & Baseline 95\% CI & \textsc{bisect}$^\dagger$ & Enacted & \textsc{bisect}$^\dagger$ & Enacted \\
\midrule
NC ($k=14$, $n=13$) & $[0, 3]$ & 4 & 2 & 6 & 2 \\
WI ($k=8$, $n=8$)   & $[0, 2]$ & 2 & 1 & 3 & 1 \\
TX ($k=38$, $n=36$) & $[1, 8]$ & 7 & 3 & 11 & 4 \\
\bottomrule
\end{tabular}
\begin{tablenotes}
\small
\item[$\dagger$] Single-run result (seed\,$=42$).
\item Strict: no incumbent home tract in district. Broad: strict open seat or incumbent
  in competitive district ($|v_D - 0.50| \leq 0.05$).
\end{tablenotes}
\end{threeparttable}
\end{table}

\paragraph{North Carolina.}
The \textsc{bisect} NC map$^\dagger$ produces 4 strict open seats and 6 broad open
seats. The two additional broad open seats are districts that contain one incumbent's
home tract but whose partisan lean shifted to competitive ($v_D \in [0.45, 0.55]$)
relative to the prior district. The enacted NC map produces only 2 strict open seats,
both of which correspond to districts where incumbents chose not to seek reelection
for personal reasons unrelated to the map.

The \textsc{bisect} strict count of 4 is above the 95\% confidence interval of $[0, 3]$
under the random baseline. This slight excess reflects the specific geography of NC:
two pairs of incumbents whose home tracts fall on the same side of the compact
geographic cuts produce pairs that leave two districts with no incumbents. This is
a known consequence of the pairing effect documented in I.1: each pairing event
creates at least one open district (the district that lost its incumbent to the pair's
district).

\paragraph{Wisconsin.}
The \textsc{bisect} WI map$^\dagger$ produces 2 strict open seats, within the $[0, 2]$
baseline range. The enacted WI map produces 1 strict open seat. The difference is modest
but directionally consistent: the algorithmic map creates twice as many strict open seats
as the enacted map.

\paragraph{Texas.}
The \textsc{bisect} TX map$^\dagger$ produces 7 strict open seats and 11 broad open
seats. The 7 strict open seats include the 4 districts paired with incumbents' home tracts
(generating 4 open districts from the 4 pairings) plus 3 additional open districts where
no incumbent happened to live. The enacted TX map produces 3 strict open seats (2 from the
new TX seats created by reapportionment, 1 from a natural retirement).

The broad open-seat count of 11 for \textsc{bisect} reflects TX's large delegation and
the significant partisan composition shifts that compact redistricting produces in the
competitive Houston, Dallas-Fort Worth, and Austin suburban rings.

\subsection{Retirement Incentive Analysis}

Table~\ref{tab:retirement} presents the estimated retirement incentive scores for each
state's incumbents under \textsc{bisect} and enacted plans. We report the mean score
across all $n$ incumbents seeking reelection and the fraction with scores above the
median threshold for historical retirement ($\text{RI} \geq 0.30$ based on calibration
against the 2002 and 2012 post-redistricting cycles).

\begin{table}[ht]
\centering
\begin{threeparttable}
\caption{Retirement incentive scores under \textsc{bisect} vs.\ enacted plans.}
\label{tab:retirement}
\begin{tabular}{l r r r r}
\toprule
 & \multicolumn{2}{c}{Mean RI score} & \multicolumn{2}{c}{Fraction RI $\geq$ 0.30} \\
\cmidrule(lr){2-3} \cmidrule(lr){4-5}
State & \textsc{bisect}$^\dagger$ & Enacted & \textsc{bisect}$^\dagger$ & Enacted \\
\midrule
NC & 0.38 & 0.12 & 0.54 & 0.15 \\
WI & 0.29 & 0.08 & 0.38 & 0.13 \\
TX & 0.31 & 0.11 & 0.42 & 0.17 \\
\bottomrule
\end{tabular}
\begin{tablenotes}
\small
\item[$\dagger$] Single-run result (seed\,$=42$).
\item RI parameters: $\alpha = 0.40$ (pairing indicator), $\beta = 0.50$ (lean change magnitude).
\end{tablenotes}
\end{threeparttable}
\end{table}

Under \textsc{bisect} maps, the mean retirement incentive score is 2.5--4.5 times
higher than under enacted maps across all three states. The fraction of incumbents
with scores above the historical retirement threshold ranges from 0.38 (WI) to 0.54
(NC) under the algorithmic plan, compared to 0.08--0.17 under enacted plans.

These scores are illustrative rather than definitive: the $\alpha$ and $\beta$
parameters are calibrated on historical averages and may not capture the specific
strategic responses of NC, WI, and TX incumbents. But the directional conclusion is
robust: algorithmic redistricting creates substantially more retirement pressure than
enacted redistricting, because it does not engineer protective districts for each
sitting member.
